\section{Conclusion}
\label{sec:conclusion}

Finite common-map compatibility and finite-data identification are different scientific achievements.  In the fixed $N=8$, $R2_{07}$ system studied here, an explicit physical CPTP map fits four sampled transitions below tolerance, yet that exact frozen map fails two genuinely unseen transitions by rigorous lower bounds above $0.62$.  The failure does not propagate to the full model class: after the new transitions are incorporated, a different whole-algebra physical map fits all six, with
\[
0.002306082156934015\ldots
\le\eps_*^{(6)}
\le0.004582638648933954\ldots<10^{-2}.
\]

The K=4 feasible set therefore exhibits finite-data non-identification, or set-membership ambiguity, at the declared tolerance.  It contains at least two distinct admissible maps with sharply different certified prospective errors; adding the pristine constraints excludes $\Phi_4$ without making the feasible set empty.  This is a finite set-membership statement, not a universal theorem about all estimators or all experimental designs.

Nothing in the finite result licenses stronger claims about map-class misspecification, arbitrary-time autonomy, semigroup structure, Markovianity or non-Markovianity, hidden state, cross-size persistence or continuum physics.  The next questions---prospective stability of a later selected witness and quantitative contraction of nested feasible sets---remain separate future studies rather than conclusions of this manuscript.
